\section{Arquitetura}
\label{sec:arquitetura}

A arquitetura separa interface, decisão, integração externa e acionamento. A
PWA não acessa SENATRAN ou dispositivo diretamente. O backend consulta um
\textit{VehicleDataProvider}, entrega apenas categorias normalizadas ao motor de
política e, se autorizado, envia ao \textit{DispenserGateway} um comando sem
placa ou conteúdo da resposta veicular. A Figura~\ref{fig:arquitetura} apresenta
essa organização.

\begin{figure}[H]
  \centering
  \noindent\resizebox{0.98\textwidth}{!}{\begin{tikzpicture}[
    node distance=6mm and 8mm,
    bloco/.style={draw, rounded corners, align=center, minimum height=9mm,
      text width=25mm, fill=black!3, font=\small},
    externo/.style={bloco, dashed},
    seta/.style={-{Latex[length=2mm]}, thick}
  ]
    \node[bloco] (ui) {PWA\\frentista/\\supervisor};
    \node[bloco, right=of ui] (api) {API e controle\\de autorização};
    \node[bloco, right=of api] (policy) {Política de\\compatibilidade};
    \node[bloco, below=of api] (audit) {Auditoria\\minimizada};
    \node[externo, above=of policy] (provider) {Provedor veicular\\simulado/\\autorizado};
    \node[bloco, right=of policy] (gateway) {Gateway do\\dispensador};
    \node[externo, right=of gateway] (device) {Simulador ou\\ESP32 de bancada};

    \draw[seta] (ui) -- (api);
    \draw[seta] (api) -- (provider);
    \draw[seta] (provider) -- (policy);
    \draw[seta] (api) -- (policy);
    \draw[seta] (policy) -- (gateway);
    \draw[seta] (gateway) -- (device);
    \draw[seta] (api) -- (audit);
    \draw[seta] (policy) -- (audit);
  \end{tikzpicture}}
  \caption{Arquitetura lógica-alvo do protótipo}
  \label{fig:arquitetura}
  \fontepropria
\end{figure}

\subsection{Interfaces principais}

A operação principal já é exposta pelo backend simulado como
\texttt{POST} \path{/api/v1/authorizations}, recebendo placa, combustível solicitado,
identificador do frentista, bomba e um UUID \path{client_request_id}. Repetir a
mesma chave com o mesmo corpo retorna a decisão persistida sem novo comando; a
reutilização com conteúdo diferente retorna conflito. A resposta tem estado,
identificador da autorização, categorias compatíveis, código de motivo e
expiração. Dados pessoais não são devolvidos à interface.

O override é uma operação separada vinculada a uma autorização em
\path{REVIEW_REQUIRED} causada por incerteza da consulta veicular. O contrato
atual exige token Bearer, identificador declarado de supervisor e justificativa;
autenticação individual permanece requisito não implementado. O contrato-alvo
do dispositivo aceitará somente
\path{LOCK_ALL} ou \path{UNLOCK(nozzle_id,ttl)}, acompanhados de
\path{command_id}, sequência monotônica, sessão, horário de emissão e
expiração, conforme o tipo da operação. Cada comando produzirá confirmação
associada ao mesmo identificador e à mesma sequência.

\subsection{Decisões tecnológicas}

O incremento atual utiliza Python/FastAPI no backend, React na PWA e SQLite para
persistência local minimizada. Um gateway em memória representa o dispensador.
Há firmware experimental que implementa o contrato MQTT para ESP32, mas ele
ainda não está integrado ao backend. PostgreSQL e contêineres permanecem opções
para uma etapa posterior, não componentes já empregados. Essas escolhas são
substituíveis: as regras de compatibilidade permanecem no domínio, e integrações
são isoladas pelas interfaces \path{VehicleDataProvider} e
\path{DispenserGateway}.

Somente o modo sintético está habilitado: a configuração do backend rejeita a
inicialização em modo SENATRAN, e o adaptador correspondente não realiza chamada
externa. Essa barreira evita uso acidental de dados ou credenciais antes da
autorização formal.

O ESP32 é o controlador inicialmente recomendado para a bancada, acionando LEDs
e, após revisão do orientador, servo ou solenoide de baixa tensão. Um Raspberry
Pi poderá atuar como gateway local, mas não como mecanismo único de segurança.
Para campo, a alternativa prioritária a estudar é integração com controlador de
pista ou dispenser certificado. A documentação do Fuel Point PLUS, por exemplo,
descreve identificação veículo-bico, consulta a servidor de autorização e
abertura da bomba somente depois da aprovação \cite{gilbarcoFuelPoint2018}; ela
serve como referência de arquitetura comercial, não como validação do protótipo.

\subsection{Disponibilidade e falhas}

O sistema não adotará \textit{fail-open}. Se o provedor estiver indisponível ou
se a resposta não puder ser mapeada de maneira inequívoca, a decisão automática
será \path{REVIEW_REQUIRED}. A organização deverá definir um procedimento
manual seguro para continuidade operacional; o protótipo apenas registrará o
override autorizado. Perda de comunicação depois da liberação fará o
controlador bloquear localmente após o TTL, independentemente do backend.

Se a PWA perder a resposta de criação, ela não presume bloqueio e não cria outra
solicitação automaticamente: a reconciliação manual repete exatamente a mesma
chave e o mesmo corpo. Se a resposta de override ou conclusão se perder, a
interface consulta o estado pelo identificador já conhecido, sem repetir o
comando físico. \path{PENDING_DEVICE} é tratado como estado físico desconhecido.
